图ANN通过有希望的邻域导航替代列表扫描，以避免评分大多数向量。
GPU难点在于“发现邻居”与“选择扩展哪个邻域”之间的数据依赖。
SONG把搜索拆成阶段来暴露距离并行，并控制动态内存开销\cite{song}。
这是替代访问策略，不是JHQ固定路由全集的另一种实现。

CAGRA针对GPU并行设计建图与搜索，覆盖小batch和大batch\cite{cagra}。
图连通性与search width形成不同于IVF probe的调参面。
边与向量的存储形成容量取舍；量化向量CAGRA同样重要，所以只测fp32图索引不足以
评价压缩索引。E1包含二者，在共同边界下报告不同工作区间，不拿跨论文提速比相除。

BANG针对另一容量区间：GPU保存压缩向量，CPU/GPU合作访问主机图\cite{bang}。
拆分工作结构获得容量，同时引入传输与协调成本。其十亿规模说明完整常驻之外的
替代方案，不构成本论文的容量证据。本文JHQ压缩搜索索引保留在单卡，
没有相同的主机图执行路径。

因此比较取决于访问规则性、表示与驻留方式，而不是图和量化的普遍优劣。
少量查询时图导航可能有吸引力，列表扫描则暴露规则编码访问与batch并行。
二者都可能受选择或移动限制；目标recall下谁更快，必须指定设备、显存和工作负载实测。
E1只测试batch 1,000，JHQ内部小batch收益不能回答相对图基线的问题。
